% from ega0-20-fr.tex
\oldpage[0\textsubscript{IV-1}]{148}
est \emph{surjectif} \sref{0.19.1.5-fr} ; comme
$\operatorname{Hom.\ cont}_C(\Omega^1_{B/A}\otimes_B C,L)=
\operatorname{Hom.\ cont}_B(\Omega^1_{B/A},L)$ par définition de la topologie
produit tensoriel, il revient au même, en vertu de \sref{0.20.4.8-fr}, de dire
que l'homomorphisme canonique
\[
  \operatorname{D\acute{e}r.\ cont}_A(C,L)\to
  \operatorname{D\acute{e}r.\ cont}_A(B,L)
\]
est surjectif. Mais la suite exacte \sref{0.20.3.8.2-fr} montre que cette
condition équivaut à $\operatorname{Exalcotop}_{B/A}(C,L)=0$, c'est-à-dire
précisément au fait que $C$ est formellement lisse relativement à $A$
\sref{0.19.9.8-fr}.
\end{proof}

\begin{corollary}[20.7.3]
\label{0.20.7.3-fr}
Supposons que dans $B$ et dans $C$, le carré d'un idéal ouvert soit ouvert.
Pour que $C$ soit une $B$-algèbre formellement lisse relativement à $A$, il
faut et il suffit que si l'on désigne par $(\mathfrak{R}_\lambda)$ un système
fondamental de voisinages de $0$ formé d'idéaux de $C$, alors, pour tout
$\lambda$, l'homomorphisme
\begin{equation}
\tag{20.7.3.1}
\label{0.20.7.3.1-fr}
  v_{C/B/A}\otimes 1_{C/\mathfrak{R}_\lambda}:
  \Omega^1_{B/A}\otimes_B(C/\mathfrak{R}_\lambda)\to
  \Omega^1_{C/A}\otimes_C(C/\mathfrak{R}_\lambda)
\end{equation}
soit inversible à gauche.
\end{corollary}

\begin{proof}
On sait en effet alors que la topologie de $\Omega^1_{B/A}$ (resp.\
$\Omega^1_{C/A}$) est déduite de celle de $B$ (resp.\ de $C$)
\sref{0.20.4.5-fr} ; on en conclut aussitôt que la topologie de
$\Omega^1_{B/A}\otimes_B C$ est aussi déduite de celle de $C$ ; le corollaire
résulte alors de \sref{0.20.7.2-fr} et \sref{0.19.1.7-fr}.
\end{proof}

\begin{proposition}[20.7.4]
\label{0.20.7.4-fr}
Soient $A$ un anneau topologique, $B$ une $A$-algèbre topologique. Pour que
$B$ soit formellement non ramifiée \sref{0.19.10.2-fr}, il faut et il suffit
que le séparé complété $\widehat{\Omega}^1_{B/A}$ soit nul.
\end{proposition}

\begin{proof}
En effet, il résulte aussitôt de \sref{0.19.10.4-fr} et
\sref{0.20.1.1-fr} que, pour que $B$ soit formellement non ramifiée, il faut et
il suffit que pour tout idéal ouvert $\mathfrak{R}$ de $B$, tout idéal ouvert
$\mathfrak{H}$ de $A$ tel que $\mathfrak{H}B\subset\mathfrak{R}$ et tout
$(B/\mathfrak{R})$-module $L$, on ait
$\operatorname{D\acute{e}r}_{A/\mathfrak{H}}(B/\mathfrak{R},L)=0$,
c'est-à-dire $\operatorname{D\acute{e}r.\ cont}_A(B,L)=0$ pour tout
$B$-module discret $L$ annulé par un idéal ouvert de $B$ ; en vertu de
\sref{0.20.4.8-fr}, cela équivaut à
$\operatorname{Hom.\ cont}_B(\Omega^1_{B/A},L)=0$ pour un tel $B$-module,
d'où aussitôt la proposition.

Lorsque $B$ est \emph{discret}, la condition de l'énoncé de
\sref{0.20.7.4-fr} équivaut donc à $\Omega^1_{B/A}=0$.
\end{proof}

\begin{corollary}[20.7.5]
\label{0.20.7.5-fr}
Soient $A$ un anneau, $\mathfrak{m}$ un idéal de $A$, $B$ une $A$-algèbre ;
munissons $A$ de la topologie $\mathfrak{m}$-préadique, $B$ de la topologie
déduite de celle de $A$, et posons $A_0=A/\mathfrak{m}$,
$B_0=B/\mathfrak{m}B=B\otimes_A A_0$. Alors, pour que $B$ soit une
$A$-algèbre formellement non ramifiée, il faut et il suffit que
$\Omega^1_{B_0/A_0}=0$ (ou encore que $B_0$ soit une $A_0$-algèbre
formellement non ramifiée).
\end{corollary}

\begin{proof}
En effet, on a
$\Omega^1_{B/A}\otimes_B B_0=\Omega^1_{B/A}/\mathfrak{m}.\Omega^1_{B/A}
=\Omega^1_{B_0/A_0}$ en vertu de \sref{0.20.5.5-fr} ; écrire que tout
sous-module ouvert de $\Omega^1_{B/A}$ est égal à $\Omega^1_{B/A}$ équivaut
par \sref{0.20.4.5-fr} à écrire que
$\mathfrak{m}.\Omega^1_{B/A}=\Omega^1_{B/A}$, d'où la conclusion.
\end{proof}

\begin{proposition}[20.7.6]
\label{0.20.7.6-fr}
Soient $u:A\to B$, $v:B\to C$ deux homomorphismes continus d'anneaux
topologiques, et supposons que $v$ fasse de $C$ une $B$-algèbre formellement
étale ; alors $v_{C/B/A}:\Omega^1_{B/A}\otimes_B C\to\Omega^1_{C/A}$ est un
bimorphisme formel \sref{0.19.1.2-fr}.
\end{proposition}

\begin{proof}
En effet, pour tout $C$-module discret $L$ annulé par un idéal ouvert de $C$,
on a alors \sref{0.20.7.4-fr}
$\operatorname{D\acute{e}r.\ cont}_B(C,L)=0$ et par suite
\sref{0.20.3.6.1-fr} l'homomorphisme canonique
